\magnification=\magstep1
\input notesmac.tex
\nopagenumbers
\centerline{\titlefont Tech and Tools, Lab 7}
\bigskip
\line{Dr.~WANG Jian-Sheng \hfill 19/21 October 1999}
\bigskip

{\sl The following problems are to be programmed in Fortran 90.}

\bigskip
%--------------------------------------------------------------------
\problem {\sc Stack}
\item{} (1) Define a data structure called stack using linked list.
Each item on the stack will be an integer.  The items are added (pushed) 
only at the top, and the last item added is removed (popped) first.
\item{} (2) Define a subroutine {\tt PUSH} that puts an integer on the stack,
and a subroutine {\tt POP} that removes an integer from the stack
(indicating error if stack is empty).
\item{} (3) Define subroutines {\tt ADD}, {\tt SUB}, {\tt MUL}, and {\tt DIV}
which take the top two integers from the stack (and also remove them
from the stack), then add, subtract, multiply, and divide,
respectively, the two numbers.  The subroutines also push the result on 
top of the stack.  The PostScript language handles calculations 
this way.
\item{} (4) Write a main program to test your subroutines.  In the main program,
first ask the user to enter arbitrary number of integers. Then arbitrary
sequence of operations ({\tt +, -, *, /}) is operated on the stack,
and the resulting stack is printed.

%--------------------------------------------------------------------
\bigskip
\problem {\sc Module}
\item{} Write a module {\tt VECTOR} for vector algebra which contains the
following:
{\parindent = 12pt
\item{} 1) Type definition of three dimensional vector $\bf a$.
\item{} 2) Define a function and overload operator {\tt +} for vector addition, $\bf a + b$.
\item{} 3) Define a function and overload operator {\tt *} for scalar multiplying a vector, $ c\; \bf a$.
\item{} 4) Define a function and overload operator {\tt *} for dot product of two vectors, $\bf a \cdot b$.
\item{} 5) Define a function and define operator {\tt .CROSS.} for cross (vector) product
of two vectors, $\bf a \times b$.
\item{} 6) Define and overload assignment {\tt =} for converting an array (of three
elements) to vector.
\item{} 7) Define and overload assignment {\tt =} for assigning all components of a
vector with the same real number, so that $\bf a = 0$ make sense. 
\item{} 8) Write a test program to check out your {\tt VECTOR} module.
}

\bye

